\section{Caso de estudio y corpus documental}

Este capítulo delimita el caso de estudio y el conjunto de materiales que sostienen el diagnóstico. Su función es acotar la investigación antes de presentar los resultados: no se busca describir la totalidad de la vivienda hñähñu del Valle del Mezquital, sino estudiar un conjunto concreto de construcciones tradicionales que todavía permiten reconocer sistemas, materiales, transformaciones y relatos asociados al saber constructivo local.

\subsection{Comunidades de estudio}

La investigación se centra en dos comunidades rurales del municipio de El Cardonal, Hidalgo: El Deca y El Buena. Ambas forman parte del contexto hñähñu del Valle del Mezquital y comparten un entorno semidesértico donde los recursos disponibles (tierra, piedra, maguey, lechuguilla, madera y cactáceas) han condicionado históricamente las formas de construir.

La selección de estas comunidades responde a tres criterios: la presencia de construcciones tradicionales aún reconocibles, la disponibilidad de registros fotográficos y geográficos, y la existencia de relatos de habitantes que permiten reconstruir usos, cambios y valoraciones locales. Esta delimitación permite tratar el caso con mayor profundidad, sin generalizar más allá de la evidencia documentada.

\subsection{Unidad de análisis}

La unidad de análisis son \emph{construcciones tradicionales} y no viviendas completas. Esta precisión es importante porque, en varios casos, lo que permanece en pie corresponde a espacios complementarios de la vida cotidiana, como cocinas, bodegas, corrales o estructuras de resguardo, más que a casas habitadas en su totalidad.

El corpus empírico inicial está conformado por cuatro construcciones documentadas mediante fotografías, ubicación geográfica y al menos un relato hablado asociado. Cada construcción será identificada mediante una clave anónima para proteger la privacidad de los propietarios y evitar la publicación de ubicaciones exactas en la versión final del documento.

\subsection{Corpus inicial de construcciones}

El corpus inicial se organizará mediante cuatro fichas base. Cada ficha corresponde a una construcción y no a una familia o tipología completa. Esta decisión permite mantener visible la evidencia de campo y evita que la clasificación se separe de los casos que la originan.

\begin{table}[ht]
	\centering
	\caption{Estructura inicial del corpus documental}
	\vspace{0.25cm}
	\label{tab:corpus_inicial}
	\begin{tabular}{@{}L{0.16\linewidth}L{0.22\linewidth}L{0.28\linewidth}L{0.24\linewidth}@{}}
		\toprule
		\textbf{Clave} & \textbf{Comunidad} & \textbf{Sistema preliminar} & \textbf{Evidencia disponible} \\
		\midrule
		C-01 & Por definir & Penca de maguey & Fotografía, ubicación y relato \\
		C-02 & Por definir & Piedra local o ``tepetate'' & Fotografía, ubicación y relato \\
		C-03 & Por definir & Bajareque o ``pajareque'' & Fotografía, ubicación y relato \\
		C-04 & Por definir & Órganos & Fotografía, ubicación y relato \\
		\bottomrule
	\end{tabular}
\end{table}

Las claves son provisionales y podrán ajustarse cuando se ordene el archivo fotográfico definitivo. En la versión pública del producto, estas claves permitirán consultar los casos sin exponer nombres de propietarios ni coordenadas exactas.

\subsection{Sistemas constructivos preliminares}

A partir de la observación de campo y de los relatos disponibles, se consideran cuatro sistemas constructivos preliminares:

\begin{enumerate}
	\item \textbf{Penca de maguey}: muros con entramado de quiote de lechuguilla y recubrimiento de pencas de maguey.
	\item \textbf{Piedra local o ``tepetate''}: muros colocados a hueso con piedra nombrada localmente como tepetate, pendiente de caracterización material precisa.
	\item \textbf{Bajareque o ``pajareque''}: entramado de lechuguilla, relleno con material vegetal o tierra, y recubrimiento térreo.
	\item \textbf{Órganos}: cerramientos vivos o semivivos formados con cactus órgano, utilizados para delimitar y proteger espacios.
\end{enumerate}

Esta lista no funciona como una clasificación cerrada. En el análisis se verificará qué rasgos materiales, constructivos y de uso permiten distinguir cada sistema y qué combinaciones aparecen entre ellos.

\subsection{Materiales de análisis}

El estudio se apoya en cuatro tipos de evidencia:

\begin{itemize}
	\item \textbf{Fotografías}: registro visual de materiales, ensambles, estado físico, cubiertas, cambios y elementos de uso.
	\item \textbf{Geolocalización}: puntos y recorridos trabajados en cartografía para ubicar las construcciones en relación con las comunidades y el territorio.
	\item \textbf{Relatos hablados}: información proporcionada por habitantes sobre usos, cambios, mantenimiento, materiales y memoria de las construcciones.
	\item \textbf{Revisión documental}: literatura sobre arquitectura vernácula, saberes constructivos, patrimonio biocultural y transformación de la vivienda rural.
\end{itemize}

La combinación de estos materiales permite construir un diagnóstico descriptivo y cartográfico, cercano al tipo de productos recurrentes en otras ICR de conservación: fichas, mapas, tipologías y criterios derivados del caso.
